% Chapter 6 -- Creating and Changing a Document (Part I, Learning; AN-14
% task form).  Sources (Law 1): os/apps/itext/itext.c (InitechText: key():
% the arrows move one character or one line, Backspace erases the
% selection or the character before the insertion point, Enter ends a
% line, Ctrl chords go only to the menus; mouseDown: click, Shift-click and
% drag selection; replace_sel: typing replaces the selection; copy_sel /
% paste: the desk Scrap; guarded(): Close and Quit ask "before quitting?",
% the close box, New and Open... ask "before closing?"; save / save_as:
% Standard File; doc_load: the saved text replaces the window's, the title
% becomes the file name; show_caret: the window follows the insertion
% point, the scroll bar is given the document height); os/flair/textedit.c
% (word wrap; FLAIR_TE_TEXT_MAX); os/flair/scrap.c (the Clipboard, 8,192
% characters, kept in memory only); os/milton/kmain.c sf_run (the Save
% dialog box, the Replace question).  Captures: screens/d_*.mouse (locked
% traces of the running desktop).  Model: Macintosh User's Guide 1991 Ch. 4
% "Creating a Document and Saving Your Work" and Ch. 5 "Changing a
% Document" (apple-macintosh.md Sec 8.2), task form of Sec 5.4.
% GATES: test-flair-itext test-itext-route test-flair-itext-scroll test-stdfile test-scrap test-pubs-flair-screens
\Gates{test-flair-itext, test-itext-route, test-flair-itext-scroll,
  test-stdfile, test-scrap, test-pubs-flair-screens}
\chapter{Creating and Changing a Document}
\label{ch:documents}

In this chapter you write a short memo with InitechText, correct it,
copy part of it to another place, save it, keep a second copy under
another name, and go back to the version you saved.  What you learn here
works the same way in every FLAIR application that works with text.

\section{Before you begin}
Start FLAIR, and close any application from the disk that is open.  The
steps begin at the desktop, with the \texttt{INITECH} icon in its usual
place at the top right of the screen.  Chapter~\ref{ch:applications}
showed you how to open InitechText; this chapter starts from there.

\section{Starting a new document}
\index{documents!new}\index{typing}
\begin{TaskSteps}
\Step Open the \file{APPS} folder, click its zoom box, and double-click
\file{TEXT.EXE}.
\Result{InitechText opens with a new document called \texttt{untitled}.
The insertion point, a thin vertical line, waits at the top left of the
window.}
\Step Type \type{To: All staff} and press \key{Enter}.
\Result{The text appears as you type, and the insertion point moves to
the start of the next line.}
\Step Type the rest of the memo, pressing \key{Enter} at the end of each
line, and twice after the second line to leave a blank line:
\begin{TypeIn}
Re: TPS reports

Every TPS report now goes out with
the new cover sheet.
\end{TypeIn}
\end{TaskSteps}

\FlairScreen[label=fig:dtyped]{d_typed}{The memo, typed}

When a line becomes too long for the window, InitechText carries the word
you are typing to the start of the next line by itself.  Press \key{Enter}
only to end a paragraph or to leave a blank line.

If you make a mistake while you type, press \key{Backspace}: it erases the
character just before the insertion point.

\section{Moving the insertion point}
\index{insertion point}
Whatever you type goes in at the insertion point.  To type somewhere else,
move the insertion point first.
\begin{TaskSteps}
\Step Click just to the right of the end of the second line, the line that
reads \texttt{Re: TPS reports}.
\Result{The insertion point moves to the end of that line.  A click to the
right of a line's end puts the insertion point at the end of the line; a
click on a character puts it just before or after the character, whichever
is nearer.}
\Step Type a space and \type{(monthly)}.
\Result{The words are added to the line.  The text after the insertion
point, if there is any, moves along to make room.}
\end{TaskSteps}

\FlairScreen[label=fig:dinsert]{d_insert}{Words added to the second line}

You can also move the insertion point with the arrow keys: \key{Left Arrow}
and \key{Right Arrow} move it one character, \key{Up Arrow} and \key{Down
Arrow} one line.  The insertion point is drawn as a steady line; it does
not blink.

\section{Selecting text}
\index{selecting!text}
To change a word or a line, you first select it.  Selected text is shown
in white on black.
\begin{TaskSteps}
\Step Point just to the left of the first word of the last line,
\texttt{the new cover sheet.}
\Step Hold down the mouse button and drag to the right, past the end of the
line, then release the button.
\Result{The whole line is selected.}
\end{TaskSteps}

\FlairScreen[label=fig:dselect]{d_select}{The last line, selected}

There are two other ways to select text:
\begin{itemize}
\item Click where the selection is to start, hold down \key{Shift}, and
  click where it is to end.  Everything between the two places is
  selected.
\item Choose \menuitem{Select All} from the \menuname{Edit} menu, or press
  \key{Ctrl-A}, to select the whole document.
\end{itemize}
To deselect the text, click anywhere in the document.

\section{Replacing and removing text}
\index{replacing text}\index{deleting text}
With text selected:
\begin{itemize}
\item Type, and what you type takes the place of the selected text.
\item Press \key{Backspace}, or choose \menuitem{Clear} from the
  \menuname{Edit} menu, and the selected text is removed.
\end{itemize}
Neither puts anything on the Clipboard.  InitechText cannot undo a change:
\menuitem{Undo} in its \menuname{Edit} menu is dimmed.  Before you make a
large change, save the document, so that you can go back to the saved
version if you need to (see ``Going back to the saved version,'' later in
this chapter).

\section{Copying text to another place}
\index{Clipboard!copying}\index{Copy command}\index{Paste command}
The \menuitem{Copy}, \menuitem{Cut} and \menuitem{Paste} commands of the
\menuname{Edit} menu let you use the same words twice, or move them, without
typing them again.  They work through the \term{Clipboard}, a holding place
in the computer's memory.
\begin{TaskSteps}
\Step Select the last line of the memo, as you did before.
\Step Choose \menuitem{Copy} from the \menuname{Edit} menu, or press
\key{Ctrl-C}.
\Result{A copy of the selected text is put on the Clipboard.  Nothing
changes on the screen.}
\Step Click below the last line of text.
\Result{The insertion point goes to the end of the document.}
\Step Press \key{Enter}, and type \type{Please read the memo about} and a
space.
\Step Choose \menuitem{Paste} from the \menuname{Edit} menu, or press
\key{Ctrl-V}.
\Result{The words on the Clipboard are put in at the insertion point.}
\end{TaskSteps}

\FlairScreen[label=fig:dpaste]{d_paste}{The copied words, pasted at the end
of a new line}

To move text instead of copying it, choose \menuitem{Cut} (\key{Ctrl-X})
in place of \menuitem{Copy}: the selected text is removed from the document
and put on the Clipboard.  Then paste it where you want it.

The Clipboard holds one piece of text at a time.  What you cut or copy
replaces what was there, and stays there, ready to be pasted again, until
you cut or copy something else.  It is kept when you begin a new document
or quit InitechText, and emptied when you restart the computer.

\section{Saving your work}
\index{saving a document}
\begin{TaskSteps}
\Step Choose \menuitem{Save} from the \menuname{File} menu, or press
\key{Ctrl-S}.
\Result{The Save dialog box asks for a name, as Chapter~\ref{ch:applications}
describes.}
\Step Erase \texttt{UNTITLED} with \key{Backspace}, type \type{cover.txt}
and press \key{Enter}.
\Result{The memo is written to the disk as \file{COVER.TXT}, and the window
takes that name.}
\end{TaskSteps}
From now on, \menuitem{Save} writes the document to \file{COVER.TXT}
without asking.  Save whenever you have done a few minutes' work that you
would not want to do again.

\section{Saving a copy under another name}
\index{Save As command}
Sometimes you want to keep a document as it is and go on working on a copy:
a memo to one department, say, as the start of a memo to another.
\begin{TaskSteps}
\Step Choose \menuitem{Save As...} from the \menuname{File} menu.
\Result{The Save dialog box opens.  Its list shows the files and folders
at the top level of the disk, \file{COVER.TXT} among them.}
\Step Erase the name in the name box and type \type{cover2.txt}.
\end{TaskSteps}

\FlairScreen[label=fig:dsaveas]{d_saveas}{Save As, with a new name typed}

\begin{TaskSteps}[resume]
\Step Click \texttt{Save}, or press \key{Enter}.
\Result{The document is saved as \file{COVER2.TXT}, and the window takes
the new name.  \file{COVER.TXT} stays on the disk as it was when you last
saved it.  Changes you make now go into \file{COVER2.TXT} when you save.}
\end{TaskSteps}
If you give the name of a file that is already on the disk, FLAIR asks
whether to replace it (Figure~\ref{fig:replace}).

\section{Going back to the saved version}
\index{documents!reverting}
If you change a document and then decide you do not want the changes, open
the saved version again and do not save the changes.
\begin{TaskSteps}
\Step Choose \menuitem{Open...} from the \menuname{File} menu, or press
\key{Ctrl-O}.
\Result{Because the document has changes you have not saved, InitechText
first asks whether to save them.}
\Step Click \texttt{Don't Save}.
\Result{The Open dialog box appears.}
\Step Double-click the name of the document, \file{COVER.TXT}.
\Result{The document comes back as it was when you last saved it.}
\end{TaskSteps}

\FlairScreen[label=fig:dreopen]{d_reopen}{The saved version, opened again;
the line added after saving is gone}

\section{Closing the document}
\index{closing a document}
InitechText holds one document in one window.  Closing the window quits
InitechText.
\begin{TaskSteps}
\Step Click the close box at the left end of the title bar.
\Result{If the document has changes you have not saved, InitechText asks
whether to save them before closing.}
\end{TaskSteps}

\FlairScreen[label=fig:dcloseask]{d_closeask}{Closing a window that holds
unsaved changes}
\Operation{FLAIR-CLOSE-ASK}

\begin{TaskSteps}[resume]
\Step Click \texttt{Save} to save the changes and close, \texttt{Don't
Save} to close without them, or \texttt{Cancel} to go on working.
\end{TaskSteps}
Choosing \menuitem{Close} or \menuitem{Quit} from the \menuname{File} menu
does the same; the question then ends \texttt{before quitting?}.

\section{Working with a long document}
\index{scrolling!a document}
A document can be longer than its window.  When it is, the scroll bar at
the right of the window becomes active, and you use it to see the rest of
the document, as you use the scroll bars of a Finder window
(Chapter~\ref{ch:windows}).

As you type past the bottom of the window, InitechText moves the contents
up to keep the insertion point in view.

\FlairScreen[label=fig:dscroll]{d_scroll}{A list longer than the window,
scrolled back by one line}

Figure~\ref{fig:dscroll} shows a list of 42 items.  The scroll box is near
the bottom of the scroll bar, because the part of the document in view is
near its end.  If lines on the screen appear drawn over one another after
you type past the bottom of the window, click one of the scroll arrows:
the window is drawn again.

A document can hold up to 4,096 characters.  When a document is full, InitechText takes no more typing or
pasting.

\section{What you have learned}
\begin{itemize}
\item Text goes in at the insertion point.  Click, or use the arrow keys,
  to move it.
\item To change text, select it first: drag across it, Shift-click, or
  choose \menuitem{Select All}.  Typing replaces the selection;
  \key{Backspace} removes it.
\item \menuitem{Copy} and \menuitem{Cut} put the selection on the
  Clipboard; \menuitem{Paste} puts the Clipboard's text at the insertion
  point.
\item \menuitem{Save} keeps your work on the disk.  \menuitem{Save As...}
  keeps a copy under another name.
\item To go back to the saved version, open the document again and click
  \texttt{Don't Save}.
\end{itemize}
